% -----------------------------------------------------------
\section{God's Nature (Nature of God)}
% -----------------------------------------------------------
	\label{GODSNATURE}
	
% ----------------------
\subsection{See also}
% ----------------------

	\hyperref[GODLINESS]{Godliness},

% % ----------------------
% \subsection{1828 American Dictionary of the English Language} % *1828
% % ----------------------
	% \label{GODSNATURE-1828}

	% \begin{compactitem}
		% \item
	% \end{compactitem}

% % ----------------------
% \subsection{Oxford English Dictionary} % *OED
% % ----------------------
	% \label{GODSNATURE-OED}

	% \begin{compactitem}
		% \item[]
	% \end{compactitem}
	
% % ----------------------
% \subsection{Restoration Lexicon} % *RL *RESTORATION LEXICON
% % ----------------------
	% \label{GODSNATURE-OED}

	% \begin{compactitem}
		% \item[]
	% \end{compactitem}

% ----------------------
\subsection{Scriptural Usage} % *SCRIPTURES
% ----------------------
	\label{GODSNATURE-SCRIPTURES}

	% % -----
	% \subsubsection{Old Testament} % *OT
	% % -----
		% \label{GODSNATURE-OT}
	
		% % --
		% \paragraph{KJV}
		% % --
			% \label{GODSNATURE-OT-KJV}
	
			% \begin{tabular}{ll}
				% Total occurrences in the OT & 
			% \end{tabular}
			
			% \begin{compactdesc}
				% \item[]
			% \end{compactdesc}
			
		% % --
		% \paragraph{Hebrew Bible}
		% % --
			% \label{GODSNATURE-HEBREW}
		
			% %
			% \subparagraph{\texorpdfstring{\he{sample word}}{}} % paste actual hebrew text here
			% %
				% \label{PAGA} % whatever is in step bible without periods
			
				% \begin{compactdesc}
					% \item[HALOT , PDF ] \textbf{} % Use the 1994 PDF
						% \begin{compactitem}
							% \item[1]
						% \end{compactitem}
					% \item[BDB , PDF ] \textbf{}
						% \begin{compactitem}
							% \item[1]
						% \end{compactitem}
					% \item[DCH PDF ] \textbf{} % don't include the actual page number, they repeat from volume to volume
						% \begin{compactitem}
							% \item[1]
						% \end{compactitem}
				% \end{compactdesc}
				
				% \begin{compactdesc}
					% \item[Some OT uses] \textbf{}
						% \begin{compactdesc}
							% \item[]
						% \end{compactdesc}
				% \end{compactdesc}
			
		% % --
		% \paragraph{Septuagint}
		% % --
			% \label{GODSNATURE-LXX}
		
			% %
			% \subparagraph{\texorpdfstring{\gr{sample word}}{}}
			% %
		
	% -----
	\subsubsection{New Testament} % *NT
	% -----
		\label{GODSNATURE-NT}
	
		% --
		\paragraph{KJV}
		% --
			\label{GODSNATURE-NT-KJV}		
	
			% \begin{tabular}{ll}
				% Total occurrences in the NT & 
			% \end{tabular}
			
			\begin{compactdesc}
				\item[{\hyperref[MARK-10-8]{Mark 10:8}}]
			\end{compactdesc}
			
		% % --
		% \paragraph{Greek New Testament} % *GREEK
		% % --
			% \label{GODSNATURE-GREEK}
		
			% %
			% \subparagraph{\texorpdfstring{\gr{sample word}}{}}
			% %
				% \label{KATALLAGE}
			
				% \begin{compactdesc}
					% \item[BDAG ] \textbf{}
						% \begin{compactitem}
							% \item 
						% \end{compactitem}
					% \item[CGL , PDF ] \textbf{}
						% \begin{compactitem}
							% \item 
						% \end{compactitem}
					% \item[LSJ , PDF ] \textbf{}
						% \begin{compactitem}
							% \item 
						% \end{compactitem}
				% \end{compactdesc}
				
				% \begin{tabular}{ll}
					% Total occurrences in the NT & 
				% \end{tabular}
				
				% \begin{compactdesc}
					% \item[Some NT uses] \textbf{}
						% \begin{compactdesc}
							% \item[]
						% \end{compactdesc}
				% \end{compactdesc}
				
			% %
			% \subparagraph{TDNT: \texorpdfstring{\gr{sample word}}{}  (3:536--556)}
			% %
				% \label{KATALLAGE-TDNT}
		
	% -----
	\subsubsection{Book of Mormon} % *BOM
	% -----
		\label{GODSNATURE-BOM}
	
		% \begin{tabular}{ll}
			% Total occurrences in the BOM & 
		% \end{tabular}
		
		\begin{compactdesc}
			\item[{\hyperref[HELAMAN-13-38]{Helaman 13:38}}]
			\item[{\hyperref[MORMON-9-19]{Mormon 9:19}}]
		\end{compactdesc}
		
	% -----
	\subsubsection{Doctrine and Covenants} % *DC
	% -----
		\label{GODSNATURE-DC}
	
		% \begin{tabular}{ll}
			% Total occurrences in the DC & 
		% \end{tabular}
		
		\begin{compactdesc}
			\item[{\hyperref[DC62-6]{DC 62:6}}]
			\item[{\hyperref[DC63-34]{DC 63:34}}]
				\begin{compactitem}
					\item Thinking specifically of how Christ is in the Father's presence here.
				\end{compactitem}
			\item[{\hyperref[DC66-12]{DC 66:12}}]
				\begin{compactitem}
					\item The only place in the scriptures that explicitly says that the Father is full of ``grace and truth.''
				\end{compactitem}
			\item[{\hyperref[DC93]{DC 93}}]
			\item[{\hyperref[DC109-77]{DC 109:77}}]
			\item[{\hyperref[DC121-28]{DC 121:28}}]
			\item[{\hyperref[DC121-32]{DC 121:32}}]
			\item[{\hyperref[DC132-20]{DC 132:20}}]
		\end{compactdesc}
		
	% % -----
	% \subsubsection{Pearl of Great Price} % *PGP
	% % -----
		% \label{GODSNATURE-PGP}
	
		% \begin{tabular}{ll}
			% Total occurrences in the PGP & 1
		% \end{tabular}
		
		% \begin{compactdesc}
			% \item[]
		% \end{compactdesc}
		
% % ----------------------
% \subsection{Key Phrases} % *KEYPHRASES *PHRASES
% % ----------------------
	% \label{GODSNATURE-PHRASES}

	% \begin{compactdesc}
		% \item[] \textbf{\#times} | list
			% % \begin{compactdesc}
				% % \item[Bible anchor verses] \phantom{.}
					% % \begin{compactdesc}
						% % \item[Romans 5:5] \gr{}
					% % \end{compactdesc}
				% % \item[Commentary] ---
			% % \end{compactdesc}
	% \end{compactdesc}
	
% % ----------------------
% \subsection{Bible Dictionaries} % *DICTIONARIES
% % ----------------------
	% \label{GODSNATURE-BD}

% % ----------------------
% \subsection{Restoration Usage} % *RESTORATION
% % ----------------------
	% \label{GODSNATURE-RESTORATION}

	% % -----
	% \subsubsection{Joseph Smith} % *JS
	% % -----
		% \label{GODSNATURE-JS}
	
		% \begin{compactdesc}
			% \item[{\hyperref[KEY]{Date/Source}}]
		% \end{compactdesc}
		
	% % -----
	% \subsubsection{Temple Ordinances}
	% % -----
		% \label{GODSNATURE-TEMPLE}
	
		% \begin{compactdesc}
			% \item[{\hyperref[KEY]{Date/Source}}]
		% \end{compactdesc}
	
	% % -----
	% \subsubsection{Brigham Young} % *BY
	% % -----
		% \label{GODSNATURE-BY}
	
		% \begin{compactdesc}
			% \item[{\hyperref[KEY]{Date/Source}}]
		% \end{compactdesc}
		
	% -----
	\subsubsection{Augustine}
	% -----
		\label{GODSNATURE-AUGUSTINE}
	
		\begin{compactdesc}
			\item[DDC I:15--16] Now although he alone is thought of as the god of gods, he is also thought of by those who imagine, invoke, and worship other gods, whether in heaven or on earth, in so far as their thinking strives to reach a being than which there is nothing better or more exalted. They are, to be sure, inspired by various ideas of excellence, of which some relate to the senses, and others to the intellect, and accordingly those who are devoted to the bodily senses85 think that either the sky itself, or the brightest element that they see in the sky, or the world itself, is the god of gods. If they try to pass beyond the visible world, they envisage something bright, and in their futile imaginations represent it either as an infinite being or as one endowed with what they see as ultimate beauty; or else they think of the figure of a human body, if they value that above all else. 16. If they do not believe in a single god of gods, but rather in many gods, or gods without number, all of them having equal status, then for these too they form a mental picture which corresponds to their various ideas of bodily excellence. Those who strive to behold the nature of God through their intellect86 place him above all visible and corporeal beings, indeed above all intelligible and spiritual beings, and above all beings that are subject to change. But they all vigorously contend for the excellence of God; it is impossible to find anyone who believes that God is a thing than which there exists something better. All, then, are agreed that what they value above all other things is God.
				\begin{compactitem}
					\item This is one reason why JS's conception of God is very different; a hierarchy of Gods---they have to differ in some respect if they are to differ in glory or power or whatever. This is one point against the sort of McConkie idea that God has to be omniscient---knowledge is one of the ways that gods can differ in the hierarchy.
				\end{compactitem}
		\end{compactdesc}
	
% ----------------------
\subsection{Commentary and Studies} % *COMMENTARY *STUDIES
% ----------------------
	\label{GODSNATURE-COMMENTARY}

	% -----
	\subsubsection{Key points to consider}
	% -----
		\label{GODSNATURE-KEYS}
	
		\begin{compactitem}
			\item \hyperref[DC39-1]{DC 39:1}
				\begin{compactitem}
					\item The ``eternity to eternity'' bit. It appears elsewhere in restoration scripture but not in the Bible.
				\end{compactitem}
		\end{compactitem}
		
	% -----
	\subsubsection{The in-dwelling relationship of Christ and the Father}
	% -----
		
		See this topic \hyperref[CHRIST-indwelling]{here}.
		
	% -----
	\subsubsection{Omniscience}
	% -----
		\label{GODSNATURE-omniscience}
		
		For God's omniscience, see: 
			\begin{compactitem}
				\item \hyperref[ALMA-18-32]{Alma 18:32}
				\item \hyperref[ALMA-26-35]{Alma 26:35}
				\item \hyperref[DC38-2]{DC 38:2}
			\end{compactitem}
		
		I read an interesting idea in a paper by Kass Gorton in spring semester 2022. Here are some excerpts from that paper:
		
			\begin{quote}
				God is designed to foreknow all of His children futures; it is assumed before you knew it, God knew you where going to pick the bag of Doritos instead of the burrito. Does God’s foreknowledge take away your freedom to choose what you really wanted for lunch? There is one more factor to add, good and evil. Good and evil spin the world round and round; any given choice always comes down to the question of right or wrong. Should you eat the burrito or the bag of Doritos? Freedom to choose exist because evil fogs God’s all knowing glasses. 

				God created all things to be good; we are created from God; therefore, we are good. We are good, but there is evil in the world; the causes of evil was not given from God himself, and He does not encourage the action of evil. Saint Augustine, an ancient theologian and philosopher, argued in his book De Libero Arbitrio (On The Free Will) why God is not the man behind the evil. He said ``God foreknows everything of which He is the cause, but He is not Himself the cause of everything He foreknows. He is not the cause of evil actions, but He is their just avenger'' (Armstrong 336). Augustine believed God does know everything an individual will do, but he is not the cause of the individuals evil actions.

				This begs the question once more, if God knows all things then how do we have freedom to choose? God is the creator of all good things; therefore, He knows all things that are good, but does He foresee the evil which He did not cause? Adam and Eve could possibly prove God does foresee evil doings because He foresaw the fall of Adam and Eve. The Christian interpretation of the fall describes the event as evil; although, latter-day scripture contradicts the assumption Adam and Eve’s transgressions were done in evil; instead the scriptures describes the action as wrongful, but justified in goodness. Written in The Book of Mormon, a latter-day scripture, it states, ``If Adam had not transgressed he would not have fallen... All things which were created must have remained in the same state in which they were after they were created... Adam fell that men might be; and men are, that they might have joy'' (The Book of Mormon, 2 Ne. 2:22,25). Adam’s transgression were done in good doing not evil doing; thus, God foresaw Adam and Eve’s transgressions because they where good. If Adam and Eve eating the fruit was evil God could not have fully foreseen it because God knows everything He causes; since He is not the cause of evil He could have not foreknow their transgressions. The same can be said for every single individual created by God. Take the burrito for example, lets assume the burrito is good because it provides all 5 food groups need to fuel your body; while the bag of Doritos could be evil because it encourages the body to be sluggish and lazy. God can foresee you eating the burrito because He is the cause of good, and God foresee what He causes; on the other hand, God’s foreknowledge of the Doritos is blurred because the Doritos is the cause of evil, and God cannot foresee things He does not cause. Creating freedom to choose because you can choose between good and evil with God’s foreknowledge of the good. 

				The freedom of choice could be put into question if it where assumed God does know the evil; it is would also be assumed Adam and Eve’s transgression was truly evil. He knew you where going to choose the Doritos over the burrito; He knew you where going to give into the evil. More importantly he knew the Doritos where evil. If that where true then your freedom to choose could not exist because God already knew your decision about the burrito and Eve’s decision about the fruit...

				Just and merciful is a characteristic Augustine and the Christian faith give to God. In the Book of Mormon, Alma, a prophet of the Lord, explains, ``All his judgments are just; that he is just in all his works, and that he has all power to save every man that believe on his name'' ( The Book Of Mormon, Alma 12.15). How is God a just and merciful god if He foresees everything? God foresees that which is good but cannot foresee that which is evil; that is proved in latter-day scripture through revelation about the fall of Adam and Eve. God can foresee us eating the burrito in the future because it is good; God cannot foresee you eating the bag of Doritos because it is the cause of evil. This gives God justification to punish you because you chose the evil over God’s foreseen good. For example, eating the bag of Doritos may lead to you being less productive throughout your day, due to the fact you did not get the proper nutrition at lunch. God is able to reward and punish because his foreknowledge of evil is blurred and his foreknowledge of good is clear. 

				It all comes down to a bag of Doritos and one burrito. Does God know what you will choose to eat? Does he take away your freedom of choice? There is no one simple answer to that difficult question, but God cannot take away our freedom of choice. Evil, in a sense, is the good that keeps free will from not existing. God foresees everything that He causes; God created all good thing, so he causes good thing; God did not create evil, so he cannot cause evil things. Since he cannot cause evil thing he can not foresee evil. God can foresee us choosing the burrito due to it goodness, but it is difficult for Him to see us choose in the bag a Doritos; due to evil you are given agency, the power to choose right from wrong. It is up to yourself to choose the path God foresees for you.
			\end{quote}